%============================================================ 7. 고찰 및 결론 (압축판)
\section{고찰 및 결론}
\label{sec:discussion}

\subsection{결과의 해석}

두 계층의 기여가 가산적이라는 것(\ref{sec:res_g3}절)은 운용상의 자유를 뜻한다. 지휘관 모델
교체와 기동 정책 재학습을 서로 기다리지 않고 독립적으로 진행할 수 있다.

지휘관의 성능을 설명하는 것은 커버리지가 아니라 재계획 간 계획 안정성이라는
관찰(\ref{sec:res_mechanism}절)은 본 문제에 국한되지 않는다. 주기적으로 재계획하는 계층
구조에서 상위 결정의 변경은 하위 계층에 전환 비용(회수되지 않는 이동과 그물)을
발생시킨다. 상위 계층이 매 재계획마다 순간 최적을 추구하면 이 비용이 누적되어, 한 번의 재계획에서 나은
배정이 에피소드 전체로는 불리해진다. 언어모델은 매 호출이 독립적이라 이 일관성이 구조적으로 보장되지
않으므로, 직전 배정을 상태에 포함시키고 담당 무리를 방위로 식별하게 한 것이 그 대응이다.
따라서 언어모델을 계획 계층에 둘 때는 계획 품질을 국소 최적성이 아니라 시간에 걸친
일관성으로 평가해야 하며, 프롬프트와 상태 표현이 그 일관성을 지원해야 한다.

가장 큰 개선이 총 선회량(\ref{sec:res_effectsize}절)에서 나타난 것은 점수맵 표현의 성질에서
온다. 경유점이 격자에 고정되므로 연속 회귀의 미세 진동이 배제되어 궤적이 부드러워지고, 선회
감소는 실제 선박에서 연료·시간·기동 여유로 환산된다. 포획당 그물 소모의 감소 역시 유효 마스크가 방위 조건과 사거리 제약을 행동 분포에 직접
부과해 무의미한 위치의 그물벽을 애초에 배제한 결과다. 다만 이 이득은 점수맵의 다봉 표현력에서 오는 것이 아니다. 유효 마스크가 남기는 셀
안에서 점수맵은 거의 단봉이며, 이득은 분포의 모양이 아니라 마스크와 격자 양자화에서 온다.

\subsection{한계와 향후 과제}
\label{sec:limitations}

검증은 시뮬레이션에 한정되며 그물 거동·얽힘 판정·관측이 단순화되어
있다(Table~\ref{tab:simfidelity}). 특히 관측 불확실성은 배정이 딛고 선 무리 구성을 흔들므로,
점수맵 정책보다 지휘관 계층에 더 크게 작용할 것이다.

가장 좋은 지휘관이 상용 폐쇄형 모델이라는 데서 오는 제약 --- 수치의 장기 재현, 능력 임계의
시점 의존성, 판단 근거 원문의 미저장 --- 은 재현성 선언에 함께 적었다.

규모는 적 10척 대 방어정 3척으로 고정이어서 방어정이 늘어 배정 조합이 급증할 때 능력 임계가
어디로 옮겨 가는지는 측정하지 않았고, 적이 규칙 기반이라 그물벽 배치를 학습해 회피하는 적응적
적은 범위 밖이다. 학습 알고리즘 비교도 같은 픽셀 행동공간 위의 MAPPO 한
종에 그치며(\ref{sec:res_algo}절), 연속 행동공간 기준선과 구성요소별 절제는 과제로 남는다.

\subsection{결론}
\label{sec:conclusion}

적보다 느린 방어정으로 자폭 군집을 막는 2계층 체계를 제안하고, 시드 대응 $2\times2$ 설계의
720 에피소드로 두 계층의 기여를 분리 측정하였다. 제안 체계는 같은 방어 수준에서 자원과 기동
비용을 크게 줄였으며, 두 계층의 기여는 가산적이었고, 언어모델 지휘관의 이득은 능력 임계 위에서만
나타났다. 실해역 이식과 관측 불확실성 하에서의 계획 안정성은 향후 과제로 남는다.
